%!TEX root = ../../main.tex
%% 1.3 연구 질문 — M-RQ1–4 문구는 MASTER_THESIS_RESEARCH_PLAN §4 고정 문구 (변경 금지)
\section{연구 질문}
\label{sec:intro-rq}

네 단계의 개선에서 다음 네 연구 질문을 도출한다. 번호 앞의 M은 학위논문(master's thesis)의 질문임을 뜻하며, 동반 저널 원고의 J-RQ 번호와 혼용하지 않는다(저널 J-RQ1은 본 논문의 M-RQ2에 대응한다).

\begin{description}[leftmargin=3.6em,style=nextline]
  \item[M-RQ1 (교전 사례 구성, I1)] 서로 다른 시간 해상도의 공개 상태·사건 기록에서 킬 조건부 교전 사례를 어떤 근거로 구성할 수 있으며, 경계와 규모 기준의 선택은 사례 구성에 어떤 변화를 만드는가?
  \item[M-RQ2 (결과 가치 구성, I2)] 최종 경기 승패와 연결한 가치모형을 이용하여 교전 전후 전략적 가치 개선을 어떻게 정의할 수 있으며, 이 측정은 관측된 결과와 얼마나 대응하고 모형·구간 선택에 얼마나 의존하는가?
  \item[M-RQ3 (사전 정보와 학습기, I3)] 이렇게 정의한 SVI는 교전 전 공개정보로 어느 정도 예측할 수 있으며, 관측 정보의 구성과 학습기의 선택은 각각 예측 성능과 확률 품질에 어떻게 반영되는가?
  \item[M-RQ4 (예측의 의미와 적용 범위, I4)] 구성한 교전 결과 예측은 초기 경기 우세에 얼마나 의존하며, 균형 상태·경기 시간·관측 조건·패치와 지역이 달라질 때 어떤 범위까지 유지되는가?
\end{description}

각 질문에는 해석의 한계가 함께 붙는다. M-RQ1에서 사람의 독립적인 경계 주석이 없는 한 ``실제 한타 탐지 정확도''를 주장하지 않으며, 규모 기준을 바꿔 생긴 코호트 사이의 절대 성능 차이를 순수한 정의 효과라고 부르지 않는다. M-RQ2에서 $\Vhat$가 경기 승패를 잘 예측하는 것과 $\SVI$가 한타 승리의 독립적 정답인 것은 별개이며, $\dV$ 자체가 인과적 기여는 아니다. M-RQ3에서 동일 성능이나 미결의 대비는 그 정보의 부재를 증명하지 않고, 같은 정보·예산에서 얻은 결과를 모든 신경망·그래프·시계열 모델로 일반화하지 않는다. M-RQ4에서 균형 구간 $\Bforty$는 사전 정보가 없는 집단도, 본질적으로 어려운 집단도 아니며, 한 셀의 유의·비유의만으로 상태별 차이를 확정하지 않는다.
